\section{Chernoff 方法与 Hoeffding 集中界}
\label{sec:05-Probability/07-Transforms-and-Bounds/04}

你在一家民调机构工作。想估计某候选人的真实支持率 \(p\)，调查了 \(n=100\) 名独立随机抽取的选民。样本支持率是 \(\hat p = \bar X_n\)，每个 \(X_i\) 是 Bernoulli\((p)\)。

假设真实支持率 \(p=0.5\)。你向客户报告：``以高概率，样本支持率与真实值之差不会超过 \(0.1\)。''客户追问：``高概率是多高？给我一个数。''

你拿出 Chebyshev。\(\Var(X_i)=p(1-p)=0.25\)，\(\Var(\hat p)=0.25/100=0.0025\)。偏差 \(\ge 0.1\) 的概率：

\[
P(|\hat p - 0.5| \ge 0.1) \le \frac{0.25}{100 \times 0.01} = 0.25.
\]

最多保证 \(75\%\) 的概率落在 \(\pm 0.1\) 以内。你以为报告了安全边界，客户听到的是``有四分之一的可能翻车''。对于一个投入了几百万广告预算的竞选决策，这个保证形同虚设。

问题出在 Chebyshev 只用了方差，把分布的全部信息压缩成了一个二阶矩。如果原始变量有界，分布的重量不能无限扩散——独立有界变量的和应该集中得更紧。Chebyshev 的 \(1/n\) 衰减反映不了这件事。信息论有一个直觉：独立事件的联合概率是因子相乘，取对数后是加法——这和矩母函数的结构异曲同工。

你需要一种能给出指数级衰减的尾部界。代价是在 Chebyshev 的假设之外，额外要求独立性或矩母函数的有界性。

\subsection{指数 Markov：把偏差放大，再卡住它}

核心想法只有四步，但每一步都值得拆开看。

对任意 \(t>0\)，事件 \(X \ge a\) 等价于 \(e^{tX} \ge e^{ta}\)。指数函数单调递增，所以没有丢失信息——这是一个等价的改写，不是放缩。

现在 \(e^{tX}\) 是非负随机变量。用 Markov：

\[
P(X \ge a) = P(e^{tX} \ge e^{ta}) \le \frac{\E[e^{tX}]}{e^{ta}} = e^{-ta} M_X(t),
\]

其中 \(M_X(t) = \E[e^{tX}]\) 是 \(X\) 的矩母函数。注意这一步的实质：你不是对 \(X\) 本身用 Markov（那只能得到 \(\E[X]/a\)，还是线性尾部），而是先通过指数函数把 \(X\) 的非线性结构——所有高阶矩——一起打包进 \(M_X(t)\)，再用 Markov。这解释了为什么 Chernoff 方法对尾部敏感：\(M_X(t) = 1 + t\E[X] + \frac{t^2}{2}\E[X^2] + \cdots\)，它编码了所有矩。Markov 只用一阶矩，Chernoff 用的是所有阶矩的组合（通过指数函数把它们压进一个数里）。

这个不等式对每一个允许的 \(t>0\)（即 \(M_X(t)<\infty\) 的那些 \(t\)）都成立。因此你可以选最紧的那一个：

\[
P(X \ge a) \le \inf_{t>0:\,M_X(t)<\infty} e^{-ta} M_X(t) = \inf_{t>0} \exp\bigl(K_X(t) - ta\bigr),
\]

其中 \(K_X(t) = \log M_X(t)\) 是累积量生成函数，在\hyperref[sec:05-Probability/07-Transforms-and-Bounds/01]{``生成函数、矩母函数与特征函数''}中引入过。指数上的 \(K_X(t)-ta\) 是 \(t\) 的函数，你的目标是找到让这个表达式最小的 \(t\)。

为什么多个自由参数 \(t\) 就能做到指数衰减？\(t\) 控制指数放大的``倍率''。\(t\) 太小，放大不够，Markov 给的界太松；\(t\) 太大，矩母函数暴涨（高次矩的贡献被指数放大），界也松。最优 \(t\) 在两者之间取得平衡——这个平衡点往往给出了与尾事件概率同量级的指数衰减率。数学上，\(\inf_{t} [K_X(t)-ta]\) 恰好是 \(K_X\) 的 Fenchel--Legendre 共轭，这建立了大偏差理论与凸分析的深层联系。

对左尾 \(P(X \le a)\)，使用 \(t<0\)，或对 \(-X\) 应用同样的右尾公式。两种方式得到对称的结构：\(P(X \le a) \le \inf_{t<0} e^{-ta}M_X(t)\)。

\subsection{独立和：为什么乘积变成加法}

Chernoff 方法真正的爆发力在独立和中释放。设 \(S_n = X_1 + \cdots + X_n\)，变量独立。

\[
M_{S_n}(t) = \E[e^{t(X_1+\cdots+X_n)}] = \E[e^{tX_1}\cdots e^{tX_n}] = \prod_{i=1}^n M_{X_i}(t).
\]

独立性让联合矩母函数分解为因子乘积——这是独立性的关键力量。如果 \(X_i\) 只是两两不相关而非独立，\(e^{tX_i}\) 之间的相关性仍然存在，乘积的期望不等于期望的乘积，分解不成立。取对数：

\[
\log M_{S_n}(t) = \sum_{i=1}^n \log M_{X_i}(t).
\]

原来每个变量的矩母函数对数都只是加法地贡献到总和上。代入 Chernoff 界：

\[
P(S_n - \E[S_n] \ge t) \le \exp\!\left(-\sup_{s>0}\bigl[st - \sum_i \log M_{X_i-\E[X_i]}(s)\bigr]\right).
\]

这解释了为什么独立和的偏离概率倾向指数衰减：矩母函数的对数与 \(n\) 成正比（每个变量贡献大致相当的 \(\log M_{X_i}(t)\)），优化后的指数项也随之与 \(n\) 成比例——于是偏离概率大致是 \(e^{-cn}\)。

但请注意一个细节：矩母函数分解时要求 \(t\) 对所有的 \(X_i\) 都有效，即所有 \(M_{X_i}(t)\) 都存在。如果某些变量的 MGF 只在某个有限区间内存在，优化 \(t\) 时就被约束在那个交集里。重尾分布的 MGF 在除 \(t=0\) 外的所有点都发散——这种情况下 Chernoff 方法不能直接用。

如果变量不独立，乘积分解失效。某些依赖结构有对应的集中不等式（如 martingale 的 Azuma--Hoeffding 利用了条件矩母函数），但直接照搬独立版本的公式是不合法的。依赖序列的集中界需要逐个场景构造，这里没有万能模板。

\subsection{Hoeffding：只假设独立和有界}

Chernoff 框架是通用模具，具体怎么铸成可用的界依赖于对 \(M_X(t)\) 做更细致的放缩。Hoeffding 不等式是这个模具在``独立 + 有界''条件下的浇铸结果。

设 \(X_1,\ldots,X_n\) 独立，且每个 \(X_i\) 几乎必然落在 \([a_i, b_i]\) 内。令 \(S_n = \sum_i X_i\)。那么

\[
P(S_n - \E[S_n] \ge t) \le \exp\!\left(-\frac{2t^2}{\sum_i (b_i-a_i)^2}\right).
\]

右侧是 Gaussian 型的尾部——指数上是 \(-t^2\) 除以一个与样本量和界宽相关的常数。双侧版本只多了一个因子 \(2\)：

\[
P(|S_n - \E[S_n]| \ge t) \le 2\exp\!\left(-\frac{2t^2}{\sum_i (b_i-a_i)^2}\right).
\]

证明骨架就是 Chernoff 四步法。关键引理（不证细节）：若 \(Y\) 满足 \(\E[Y]=0\) 且 \(Y\in[a,b]\) a.s.，则

\[
\E[e^{tY}] \le e^{t^2(b-a)^2/8}.
\]

这个引理本身非常有味道：它说一个零均值有界变量的矩母函数不会超过一个方差为 \((b-a)^2/4\) 的 Gaussian 的矩母函数（后者是 \(e^{t^2\sigma^2/2}\)，代 \(\sigma^2=(b-a)^2/4\) 就得到指数上的 \((b-a)^2/8\)）。有界性比 Gaussian 尾巴薄，矩母函数增长自然更慢——这是直观。形式上看，Hoeffding 引理的证明依赖于指数函数的凸性和对称化（引入独立复制 \(Y'\)），借助凸组合在 \([a,b]\) 上的极值必然在端点达到来放缩。逐因子相乘后，指数中的系数从 \(\sum (b_i-a_i)^2/8\) 变成了分母中的 \(\sum (b_i-a_i)^2\)。

最常用的特例是 \(X_i \in [0,1]\)。此时 \(\sum_i (b_i-a_i)^2 = n\)。把事件 \(|\bar X_n - \E[\bar X_n]| \ge \varepsilon\) 改写为 \(|S_n - \E[S_n]| \ge n\varepsilon\)，代入：

\[
P(|\bar X_n - \E[\bar X_n]| \ge \varepsilon) \le 2e^{-2n\varepsilon^2}.
\]

数字对比最能说明问题。回到 \(n=100\)，\(\varepsilon=0.1\)，\(p=0.5\)：

\begin{itemize}
\tightlist
\item
  Chebyshev：\(\le 0.25\)；
\item
  Hoeffding：\(\le 2e^{-2\times 100\times 0.01} = 2e^{-2} \approx 0.271\)。
\end{itemize}

在 \(n=100\) 时二者差不多——指数衰减还来不及通过一次项展开与多项式的 \(1/n\) 拉开差距。

但 \(n=1000\) 时：

\begin{itemize}
\tightlist
\item
  Chebyshev：\(\Var(\bar X_n) = 0.25/1000\)，界 \(= 0.25/(1000\times 0.01) = 0.025\)；
\item
  Hoeffding：\(2e^{-2\times 1000\times 0.01} = 2e^{-20} \approx 4.1\times 10^{-9}\)。
\end{itemize}

Chebyshev 还在百分之几的级别磨蹭，Hoeffding 已经渺小到十亿分之一。\(n=10000\) 时，Chebyshev 给出 \(0.0025\)，Hoeffding 给出 \(2e^{-200}\approx 2.8\times 10^{-87}\)——天文数字级别的差距。指数衰减一开始可能不比多项式快，但 \(n\) 大起来之后是指数对多项式的绝对碾压。这就是为什么统计学习和高维概率论中，集中不等式如此的不可或缺：当你想联合控制成千上万个事件的错误概率时，指数衰减几乎是唯一出路。

\subsection{分布无关：不用知道分布也能保证样本量}

Hoeffding 不等式的一个直接工程价值是：你不需要知道 \([0,1]\) 内 Bernoulli 的具体 \(p\) 值，甚至不需要知道变量是 Bernoulli 还是其他 \([0,1]\) 内的分布。只要独立、有界这两个条件成立，界就有效。

客户追问：``到底要调查多少人，才能有 \(95\%\) 的把握让误差不超过 \(3\%\)？''

翻译成不等式：需要 \(P(|\hat p - p| \ge 0.03) \le 0.05\)。用 Hoeffding：

\[
2e^{-2n\cdot 0.03^2} \le 0.05.
\]

\[
e^{-0.0018n} \le 0.025,\qquad -0.0018n \le \log 0.025 \approx -3.689,
\]

\[
n \ge \frac{3.689}{0.0018} \approx 2050.
\]

约需调查 \(2050\) 人。这个数字仅取决于你对精度 \(\varepsilon\) 和置信度 \(\delta\) 的要求，与 \(p\) 无关——这是``分布无关''的真正含义。换成 Chebyshev 做同样计算：

\[
\frac{0.25}{n\cdot 0.03^2} \le 0.05 \;\Longrightarrow\; n \ge \frac{0.25}{0.0009\times 0.05} \approx 5556.
\]

Chebyshev 要五千多人。关键是当 \(\delta\) 很小时（比如要求 \(99\%\) 置信度，即 \(\delta=0.01\)），Chebyshev 要求的样本量 \(\propto 1/\delta = 100\)（相对基准的倍数），Hoeffding 只要求 \(\propto \log(1/\delta) \approx 4.6\)。置信度从 \(95\%\) 提到 \(99.9\%\)，Chebyshev 样本量暴增 \(10\) 倍，Hoeffding 只增约 \(1.5\) 倍。这是指数尾部对多项式尾部的结构性优势。

但``分布无关''的前提仍然是严格的。独立随机抽样、每个受访者诚实回答、目标总体在这段时间内稳定——缺一不可。如果样本存在自相关（比如在同一社区集中抽样）、选择偏差（只在线调查）或分布漂移（选举前一周舆论突变），公式给出的保证不成立。``分布无关''省去了估计方差的需要，但绝不省去核实抽样框架的责任。

\subsection{Bernoulli 和的显式 Chernoff 界}

上一节的 Binomial 例子用的是 Hoeffding（通用有界变量的放缩版）。如果精确知道变量是 Bernoulli，可以直接对 MGF 做精确优化而不借助 Hoeffding 引理的放缩，得到更紧的乘法形式。

设 \(S_n \sim \text{Binomial}(n,p)\)，\(\mu = np\)。对相对偏差 \(\delta>0\)，

\[
P(S_n \ge (1+\delta)\mu) \le \left(\frac{e^\delta}{(1+\delta)^{1+\delta}}\right)^{\!\mu}.
\]

这个界比 Hoeffding 紧，因为它利用了 Bernoulli 的精确 MGF：\(M_{X_i}(t) = 1-p+pe^t = 1 + p(e^t-1) \le e^{p(e^t-1)}\)（利用了 \(1+x \le e^x\)），优化 \(t\) 直接得到上式。右端的函数 \((\frac{e^\delta}{(1+\delta)^{1+\delta}})^\mu\) 当 \(\delta\) 较小时可以进一步放缩为更易记的形式：

\[
P(S_n \ge (1+\delta)\mu) \le e^{-\mu\delta^2/3},\qquad 0<\delta\le 1.
\]

验证：令 \(f(\delta) = \delta - (1+\delta)\log(1+\delta)\)，\(f(0)=0\)，\(f'(\delta) = -\log(1+\delta)\)，\(f''(\delta) = -1/(1+\delta) \le -1/2\)（当 \(\delta\le 1\)），两次积分得 \(f(\delta) \le -\delta^2/3\)。右边少了 \(3\) 的因子——比 Gaussian 尾部（\(e^{-t^2/2}\)）略松，但胜在形式统一且通用。

精确优化中还会出现 KL 散度。把链式推导写全：对 Bernoulli 变量，Chernoff 界的优化给出

\[
\frac{1}{n}\log P(\bar X_n \ge q) \le -D(q\|p),\qquad q>p,
\]

其中 \(D(q\|p) = q\log\frac{q}{p} + (1-q)\log\frac{1-q}{1-p}\) 是 Bernoulli 分布之间的 KL 散度。大偏差概率的指数衰减率恰好是这两个分布的信息散度——不是巧合。\hyperref[sec:07-Information-Theory/03-Divergences/01]{``散度''}中会系统展开：Sanov 定理进一步把这种联系推广到任意经验分布的大偏差，大偏差速率函数正是 KL 散度。

\subsection{一个端到端的数值比较：同一问题，不同界}

把民调问题放在不同 \(n\) 下，用四种界同时算出数值保证，感受一下差距。

问题：\(X_i \sim \text{Bernoulli}(0.5)\) i.i.d.，求 \(P(|\bar X_n - 0.5| \ge 0.1)\) 的上界。

\begin{itemize}
\tightlist
\item
  \textbf{Chebyshev}：\(\frac{0.25}{n \cdot 0.01}\)。
\item
  \textbf{Hoeffding}：\(2e^{-2n\cdot 0.01}\)。
\item
  \textbf{Binomial 精确 Chernoff}（乘法形式）：对右尾用 \((\frac{e^{0.2}}{1.2^{1.2}})^{0.5n}\)，再翻倍得双侧。
\item
  \textbf{精确值}（Binomial PMF 累加）：\(P(|\bar X_n - 0.5| \ge 0.1) = 2\sum_{k \le \lfloor 0.4n \rfloor} \binom{n}{k}0.5^n\)（检验边界）。
\end{itemize}

\(n=10\)：
\begin{itemize}
\tightlist
\item Chebyshev：\(0.25/(10\cdot 0.01) = 2.5\)（截断为 \(1\)，等于没说）；
\item Hoeffding：\(2e^{-0.2} \approx 1.637\)（同样截断为 \(1\)）；
\item 精确值：\(P(|\bar X_{10}-0.5|\ge 0.1) = P(S_{10}\le 4 \text{ 或 } \ge 6) = 2\cdot P(S_{10}\le 4) \approx 0.754\)——因为有界离散分布的实际离散性。
\end{itemize}
\(n=10\) 时，所有界都松得像废话——因为 \(n\) 太小，二项分布的质量还高度离散。

\(n=100\)：
\begin{itemize}
\tightlist
\item Chebyshev：\(0.25\)；
\item Hoeffding：\(2e^{-2} \approx 0.271\)；
\item 精确值（近似）：约 \(0.058\)（用正态近似 + 连续性修正估算）。
\end{itemize}
Chebyshev 还在四分之一级别，Hoeffding 已经跟真实概率在数量级上接近了（差约 \(5\) 倍）。

\(n=1000\)：
\begin{itemize}
\tightlist
\item Chebyshev：\(0.025\)；
\item Hoeffding：\(2e^{-20} \approx 4.1\times 10^{-9}\)；
\item 精确值：\[
P(|\bar X_{1000}-0.5|\ge 0.1) = P(|S_{1000}-500|\ge 100) \approx 2(1-\Phi(100/\sqrt{250})) \approx 2(1-\Phi(6.32)) \approx 2.6\times 10^{-10}.
\]
\end{itemize}
在这个量级上，Hoeffding 和精确值只差一个数量级左右，Chebyshev 差了几个宇宙。

这个对照表不是让你记住具体数字，而是建立直觉：在 \(n\) 较小的区域（几十到上百），Chebyshev 和 Hoeffding 的实际差距并不天壤之别——因为指数还没有完全脱离多项式区域。一旦 \(n\) 跨过几百的门槛，指数衰减开启碾压模式。这个门门槛的位置取决于变量的界宽和方差——分布越集中，指数优势越早显现。

\subsection{界的选择是一种信息预算}

没有一种界在所有场景下都是最优的。选择集中界就像花一笔信息预算：你掌握的分布属性越多，你能``买''到的尾部保证越强。但这笔预算不是免费的——更强的假设往往意味着更难验证。

\begin{itemize}
\tightlist
\item
  只知道多个坏事件 → Union Bound：\(P(\cup A_i) \le \sum P(A_i)\)。一条信息不用，界也最松，但假设最轻。
\item
  非负 + 知道均值 → Markov：\(P(X\ge a) \le \E[X]/a\)。只用一阶矩，\(O(1/a)\) 尾部。
\item
  知道均值 + 方差 → Chebyshev / Cantelli：\(O(1/a^2)\) 尾部。Cantelli 单侧版本给出 \(\sigma^2/(\sigma^2+a^2)\)，在 \(a\) 较小时比双侧 Chebyshev 紧得多。
\item
  独立 + 有界 → Hoeffding：\(O(e^{-ca^2})\) 尾部。样本量保障公式给出 \(n \ge \frac{\log(2/\delta)}{2\varepsilon^2}\)。
\item
  知道精确 MGF → 优化 Chernoff：指数衰减率由 KL 散度表达，在中等 \(n\) 下比 Hoeffding 紧。
\item
  还掌握方差的局部结构 + 更多矩信息 → Bernstein / Bennett：在某些区域（如大方差但薄尾）能同时利用方差和界宽信息，优于纯 Hoeffding。
\end{itemize}

不要默认越复杂的界越好。如果你的 \(n\) 只有 \(20\)，Chebyshev、Hoeffding 和精确 Chernoff 的数值差异在绝对意义上可能都不大——区别在于相对保守程度。如果你无法验证独立性，套用 Hoeffding 只是把自欺包装成定理。如果你的数据有明显的时间序列结构，直接用 martingale 版本的集中不等式（Azuma），或退回到对依赖结构做专门的方差估计。可靠而松的界比紧但建立在可疑假设上的界更有决策价值。

下一单元转向极限定理：大数定律和中心极限定理。它们不提供有限 \(n\) 的显式上界——那是本节的领域，Chernoff--Hoeffding 正是为有限样本提供非渐近的指数级保证——但极限定理揭示了长期平均和波动形状在 \(n\to\infty\) 时的普适规律，二者共同构成``随机量随样本量增大如何行为''的完整拼图。
